% scoring_v2_spec.tex
% Formal specification of the scoring_v2 evaluation pipeline.
% Compile with: pdflatex scoring_v2_spec.tex (twice for references)

\documentclass[11pt, a4paper]{article}

% ── Packages ────────────────────────────────────────────────────────────────
\usepackage[margin=2.5cm]{geometry}
\usepackage{amsmath, amssymb, amsthm}
\usepackage{mathtools}          % dcases, coloneqq
\usepackage{algorithm}
\usepackage{algpseudocode}
\usepackage{booktabs}
\usepackage{array}
\usepackage{xcolor}
\usepackage[colorlinks=true, linkcolor=blue!60!black, citecolor=green!50!black]{hyperref}

% ── Theorem environments ─────────────────────────────────────────────────────
\theoremstyle{definition}
\newtheorem{definition}{Definition}[section]
\newtheorem{remark}{Remark}[section]

% ── Operators ────────────────────────────────────────────────────────────────
\DeclareMathOperator*{\argmax}{arg\,max}
\newcommand{\1}[1]{\mathbf{1}\!\left[#1\right]}   % indicator
\newcommand{\Geo}{\mathcal{G}}                      % geometric mean
\newcommand{\Hit}{\mathcal{H}}                      % hit set
\newcommand{\Var}{\mathcal{V}}                      % variant set
\newcommand{\Methods}{\mathcal{M}}                  % method set
\newcommand{\Desc}{\mathcal{D}}                     % descriptor set
\newcommand{\fp}[1]{\phi(#1)}                       % fingerprint

% ── Algorithmic settings ─────────────────────────────────────────────────────
\algrenewcommand\algorithmicrequire{\textbf{Input:}}
\algrenewcommand\algorithmicensure{\textbf{Output:}}
\algnewcommand\algorithmicreturn{\textbf{return} }
\algnewcommand\Return{\State\algorithmicreturn}

% ────────────────────────────────────────────────────────────────────────────

\title{%
  \textbf{Formal Specification of the \texttt{scoring\_v2} Evaluation Pipeline}\\[0.5ex]
  \large Expanding Synthesizable Chemical Space --- CSC2541H
}
\author{Tyler Biswurm \and Jingxin Huang \and Yuma Nakamura \and Ben Luo}
\date{March 2026}

\begin{document}
\maketitle
\tableofcontents
\bigskip

% ════════════════════════════════════════════════════════════════════════════
\section{Notation and Setup}
\label{sec:notation}
% ════════════════════════════════════════════════════════════════════════════

Let $\mathcal{S}$ denote the set of all valid SMILES strings and
$\mathcal{M}$ the set of all chemically valid molecules.
For each molecule $m \in \mathcal{M}$, we write $v_d(m) \in \mathbb{R}$
for the value of physicochemical descriptor $d$.

The descriptor set used throughout this pipeline is
\[
  \Desc \;=\; \{\,\text{MW},\; \text{CLogP},\; \text{TPSA},\;
                  \text{HBD},\; \text{HBA},\; \text{RB}\,\},
\]
where MW is molecular weight (Da), CLogP is the Wildman--Crippen
partition coefficient, TPSA is topological polar surface area (\AA$^2$),
HBD and HBA are the numbers of hydrogen-bond donors and acceptors
respectively, and RB is the number of rotatable bonds.
All descriptors are computed using RDKit~\cite{rdkit}.

The six tolerance parameters $\delta_d$ corresponding to the acceptable
deviation from a property target are listed in Table~\ref{tab:tolerances}.

\begin{table}[h]
\centering
\caption{Per-descriptor tolerance parameters used in the physicochemical
         conservation score (Section~\ref{sec:desirability}).}
\label{tab:tolerances}
\begin{tabular}{llcc}
\toprule
Descriptor & Symbol & Unit & Tolerance $\delta_d$ \\
\midrule
Molecular weight           & $\text{MW}$   & Da   & 50.0 \\
Wildman--Crippen LogP      & $\text{CLogP}$& --   & 1.0  \\
Topological polar surface area & $\text{TPSA}$ & \AA$^2$ & 20.0 \\
H-bond donors              & $\text{HBD}$  & --   & 1.0  \\
H-bond acceptors           & $\text{HBA}$  & --   & 2.0  \\
Rotatable bonds            & $\text{RB}$   & --   & 2.0  \\
\bottomrule
\end{tabular}
\end{table}

The default evaluation thresholds and reporting configurations are
given in Table~\ref{tab:thresholds}.

\begin{table}[h]
\centering
\caption{Default threshold and binning parameters.}
\label{tab:thresholds}
\begin{tabular}{lll}
\toprule
Parameter & Values & Description \\
\midrule
$\tau_t$ & $\{0.6,\; 0.7,\; 0.85\}$ & Substructural conservation threshold \\
$\tau_d$ & $0.8$ & Physicochemical conservation threshold \\
Baseline quality bins & $[0,0.5),\;[0.5,0.7),\;[0.7,0.85),\;[0.85,1]$
                      & Stratification bins (Section~\ref{sec:summary}) \\
Fingerprint radius & $r = 2$ & ECFP4 Morgan radius \\
Fingerprint length & $B = 2048$ & Bit-vector dimension \\
\bottomrule
\end{tabular}
\end{table}


% ════════════════════════════════════════════════════════════════════════════
\section{Property Specification}
\label{sec:spec}
% ════════════════════════════════════════════════════════════════════════════

\begin{definition}[Property specification, \texttt{MolSpec}]
\label{def:molspec}
A \emph{property specification} $\sigma$ is a tuple
\[
  \sigma \;=\; \bigl(\mathbf{f}_\sigma,\;
                      \mu_d^\sigma\bigr)_{d \in \Desc},
\]
where $\mathbf{f}_\sigma \in \{0,1\}^B$ is the ECFP4 fingerprint of the
reference molecule (Section~\ref{sec:ecfp4}), and
$\mu_d^\sigma \in \mathbb{R}$ is the target value of descriptor $d$.
Optionally, $\sigma$ also stores a three-dimensional conformer
$\hat{m}_\sigma$ for ESPsim evaluation (Section~\ref{sec:espsim}).
\end{definition}

The reference molecule from which $\sigma$ is derived is a ChEMBL bioactive
compound selected for its demonstrated pharmacological activity.
Its molecular identity is retained only as a diagnostic reference; it is
not an optimization target.

\begin{algorithm}[H]
\caption{\textsc{MakeSpec}$(s,\, \textit{conformer})$ --- construct a property specification}
\label{alg:makespec}
\begin{algorithmic}[1]
\Require SMILES string $s \in \mathcal{S}$;
         boolean \textit{conformer} (generate 3D conformer for ESPsim)
\Ensure  Property specification $\sigma$, or $\bot$ if $s$ is invalid
\State $m \leftarrow \textsc{ParseSMILES}(s)$
\If{$m = \bot$} \Return $\bot$ \EndIf
\State $\mathbf{f}_\sigma \leftarrow \textsc{ECFP4}(m)$
       \Comment{Algorithm~\ref{alg:ecfp4}}
\For{$d \in \Desc$}
  \State $\mu_d^\sigma \leftarrow v_d(m)$
\EndFor
\State $\hat{m}_\sigma \leftarrow \textsc{ETKDG}(m)$ \textbf{if} \textit{conformer} \textbf{else} $\varnothing$
       \Comment{Algorithm~\ref{alg:etkdg}}
\Return $\sigma = \bigl(\mathbf{f}_\sigma,\, (\mu_d^\sigma)_{d \in \Desc},\, \hat{m}_\sigma\bigr)$
\end{algorithmic}
\end{algorithm}


% ════════════════════════════════════════════════════════════════════════════
\section{Molecular Fingerprint (ECFP4)}
\label{sec:ecfp4}
% ════════════════════════════════════════════════════════════════════════════

\begin{definition}[ECFP4 fingerprint]
\label{def:ecfp4}
The \emph{ECFP4 fingerprint} of molecule $m$ is a bit-vector
$\fp{m} \in \{0,1\}^B$, $B = 2048$, computed by Morgan's circular
algorithm with radius $r = 2$.
At iteration $r$, each atom $a$ accumulates a hash identifier
$h_a^{(r)}$ encoding the multi-set of identifiers in its $r$-hop
neighbourhood:
\[
  h_a^{(r)} \;=\; \textsc{Hash}\!\left(h_a^{(r-1)},\;
    \bigl\{\, h_b^{(r-1)} : b \in \mathcal{N}(a) \,\bigr\}\right),
  \qquad h_a^{(0)} = \textsc{AtomInvariant}(a),
\]
where $\mathcal{N}(a)$ denotes the set of atoms adjacent to $a$.
Each resulting identifier is mapped to a bit position via
$\textsc{Hash}(\cdot) \bmod B$, and the corresponding bit is set.
\end{definition}

\begin{algorithm}[H]
\caption{\textsc{ECFP4}$(m)$ --- compute the ECFP4 fingerprint}
\label{alg:ecfp4}
\begin{algorithmic}[1]
\Require Molecule $m \in \mathcal{M}$
\Ensure  Bit-vector $\mathbf{f} \in \{0,1\}^B$, $B = 2048$
\State $\mathbf{f} \leftarrow \mathbf{0}^B$
\For{$a \in \text{atoms}(m)$}
  \State $h_a^{(0)} \leftarrow \textsc{AtomInvariant}(a)$
\EndFor
\For{$r = 1$ \textbf{to} $2$}
  \For{$a \in \text{atoms}(m)$}
    \State $h_a^{(r)} \leftarrow \textsc{Hash}\!\left(h_a^{(r-1)},\;
           \bigl\{h_b^{(r-1)} : b \in \mathcal{N}(a)\bigr\}\right)$
    \State $\mathbf{f}\!\left[h_a^{(r)} \bmod B\right] \leftarrow 1$
  \EndFor
\EndFor
\Return $\mathbf{f}$
\end{algorithmic}
\end{algorithm}


% ════════════════════════════════════════════════════════════════════════════
\section{Substructural Conservation Score (Tanimoto)}
\label{sec:tanimoto}
% ════════════════════════════════════════════════════════════════════════════

\begin{definition}[Tanimoto similarity]
\label{def:tanimoto}
For two bit-vectors $\mathbf{a}, \mathbf{b} \in \{0,1\}^B$,
the \emph{Tanimoto (Jaccard) coefficient} is
\[
  T(\mathbf{a}, \mathbf{b})
  \;=\; \frac{\mathbf{a} \cdot \mathbf{b}}
             {\|\mathbf{a}\|_1 + \|\mathbf{b}\|_1 - \mathbf{a} \cdot \mathbf{b}}
  \;=\; \frac{|A \cap B|}{|A \cup B|},
\]
where $A = \{i : a_i = 1\}$ and $B = \{i : b_i = 1\}$ are the
support sets of $\mathbf{a}$ and $\mathbf{b}$, and
$\|\mathbf{a}\|_1 = \sum_i a_i$.
\end{definition}

\begin{remark}
$T \in [0, 1]$; $T = 1$ iff $\mathbf{a} = \mathbf{b}$;
$T = 0$ iff $A \cap B = \varnothing$.
\end{remark}

The \emph{substructural conservation score} of variant $m$ against
specification $\sigma$ is
\[
  \boxed{T(m, \sigma) \;=\; T\!\bigl(\fp{m},\, \mathbf{f}_\sigma\bigr).}
\]

\begin{algorithm}[H]
\caption{\textsc{TanimotoToSpec}$(m, \sigma)$ --- substructural conservation score}
\label{alg:tanimoto}
\begin{algorithmic}[1]
\Require Molecule $m \in \mathcal{M}$; specification $\sigma$
\Ensure  $T(m, \sigma) \in [0, 1]$
\State $\mathbf{f} \leftarrow \textsc{ECFP4}(m)$
\State $n_a \leftarrow \|\mathbf{f}\|_1$;\quad
       $n_b \leftarrow \|\mathbf{f}_\sigma\|_1$;\quad
       $n_{ab} \leftarrow \mathbf{f} \cdot \mathbf{f}_\sigma$
\Return $n_{ab} \;/\; (n_a + n_b - n_{ab})$
\end{algorithmic}
\end{algorithm}


% ════════════════════════════════════════════════════════════════════════════
\section{Physicochemical Conservation Score (Desirability)}
\label{sec:desirability}
% ════════════════════════════════════════════════════════════════════════════

\subsection{Linear tolerance function}

\begin{definition}[Linear tolerance function]
\label{def:linear-tol}
For descriptor $d$ with target $\mu \in \mathbb{R}$ and tolerance
$\delta > 0$, the \emph{linear tolerance desirability} of value $v$ is
\[
  \ell(v,\, \mu,\, \delta)
  \;=\; \max\!\left(0,\; 1 - \frac{|v - \mu|}{\delta}\right) \;\in\; [0, 1].
\]
\end{definition}

The function equals 1 when $v = \mu$, decreases linearly, and reaches 0
when $|v - \mu| \geq \delta$.
For integer-valued descriptors (HBD, HBA, RB), the same formula applies
with real-valued arguments.

\subsection{Geometric mean aggregation}

\begin{definition}[Geometric mean]
\label{def:geomean}
For a vector $(s_1, \ldots, s_k) \in [0,1]^k$, the \emph{geometric mean} is
\[
  \Geo(s_1, \ldots, s_k)
  \;=\; \begin{dcases}
    \exp\!\left(\frac{1}{k}\sum_{i=1}^k \ln s_i\right)
      & \text{if } s_i > 0 \;\text{ for all } i, \\
    0 & \text{otherwise.}
  \end{dcases}
\]
\end{definition}

\begin{remark}
The geometric mean is the appropriate aggregation for desirability functions
because a single component equal to zero drives the aggregate to zero,
reflecting the requirement that \emph{all} physicochemical properties must
be conserved simultaneously. The arithmetic mean would allow a high score on
one descriptor to compensate for a violation on another.
\end{remark}

\subsection{Composite physicochemical conservation score}

\begin{definition}[Physicochemical conservation score]
\label{def:desirability}
The \emph{physicochemical conservation score} of variant $m$ against
specification $\sigma$ is
\[
  \boxed{
    \Psi(m, \sigma)
    \;=\; \Geo\!\Bigl(\ell\!\bigl(v_d(m),\, \mu_d^\sigma,\, \delta_d\bigr)\Bigr)_{d \in \Desc}
    \;\in\; [0, 1].
  }
\]
\end{definition}

\begin{algorithm}[H]
\caption{\textsc{Desirability}$(m, \sigma)$ --- physicochemical conservation score}
\label{alg:desirability}
\begin{algorithmic}[1]
\Require Molecule $m \in \mathcal{M}$; specification $\sigma$
\Ensure  $\Psi(m, \sigma) \in [0, 1]$;
         component scores $(s_d)_{d \in \Desc} \in [0,1]^6$
\For{$d \in \Desc$}
  \State $s_d \leftarrow \ell\!\bigl(v_d(m),\; \mu_d^\sigma,\; \delta_d\bigr)$
\EndFor
\Return $\Geo(s_{\text{MW}},\, s_{\text{CLogP}},\, s_{\text{TPSA}},\,
             s_{\text{HBD}},\, s_{\text{HBA}},\, s_{\text{RB}})$
\end{algorithmic}
\end{algorithm}


% ════════════════════════════════════════════════════════════════════════════
\section{Three-Dimensional Similarity (ESPsim, Supplementary)}
\label{sec:espsim}
% ════════════════════════════════════════════════════════════════════════════

\subsection{Conformer generation}

\begin{algorithm}[H]
\caption{\textsc{ETKDG}$(m)$ --- generate a 3D conformer}
\label{alg:etkdg}
\begin{algorithmic}[1]
\Require Molecule $m \in \mathcal{M}$
\Ensure  3D conformer $\hat{m}$, or $\varnothing$ on failure
\State $m' \leftarrow \textsc{AddHydrogens}(m)$
\State Set random seed to 42
\State $\hat{m} \leftarrow \textsc{EmbedMolecule}(m',\, \text{ETKDGv3})$
\If{embedding failed} \Return $\varnothing$ \EndIf
\State $\hat{m} \leftarrow \textsc{MMFFOptimize}(\hat{m})$
\Return $\hat{m}$
\end{algorithmic}
\end{algorithm}

\subsection{ESPsim score}

\begin{definition}[ESPsim score]
\label{def:espsim}
Let $\hat{m}$ and $\hat{m}_\sigma$ be 3D conformers of variant $m$ and
reference $m_\sigma$ respectively.
The \emph{ESPsim score} is the arithmetic mean of the shape similarity
$S_\text{shape}$ and the electrostatic potential similarity $S_\text{ESP}$:
\[
  E(m, \sigma)
  \;=\; \frac{S_\text{shape}(\hat{m},\, \hat{m}_\sigma)
              + S_\text{ESP}(\hat{m},\, \hat{m}_\sigma)}{2}
  \;\in\; [0, 1].
\]
Both $S_\text{shape}$ and $S_\text{ESP}$ are computed by the
\texttt{espsim} library~\cite{espsim}.
\end{definition}

\begin{algorithm}[H]
\caption{\textsc{ESPsim}$(m, \sigma)$ --- 3D shape and electrostatic similarity}
\label{alg:espsim}
\begin{algorithmic}[1]
\Require Molecule $m \in \mathcal{M}$; specification $\sigma$ with
         $\hat{m}_\sigma \neq \varnothing$
\Ensure  $E(m, \sigma) \in [0, 1]$, or $\bot$ on failure
\If{\texttt{espsim} not installed \textbf{or} $\hat{m}_\sigma = \varnothing$}
  \Return $\bot$
\EndIf
\State $\hat{m} \leftarrow \textsc{ETKDG}(m)$
\If{$\hat{m} = \varnothing$} \Return $\bot$ \EndIf
\State $S_\text{shape} \leftarrow \textsc{GetShapeSim}(\hat{m},\, \hat{m}_\sigma)$
\State $S_\text{ESP}   \leftarrow \textsc{GetEspSim}(\hat{m},\, \hat{m}_\sigma)$
\Return $(S_\text{shape} + S_\text{ESP}) / 2$
\end{algorithmic}
\end{algorithm}

\begin{remark}
ESPsim is a supplementary diagnostic only.
It couples evaluation to the specific 3D geometry of the ChEMBL reference
conformer, which may not represent the bioactive conformation.
It is not used in hit classification.
\end{remark}


% ════════════════════════════════════════════════════════════════════════════
\section{Single-Variant Scoring}
\label{sec:score-variant}
% ════════════════════════════════════════════════════════════════════════════

The following algorithm assembles the full feature vector for a single
variant.
Synthesizability (Gate~1) is the caller's responsibility: only variants that
have already passed AiZynthFinder retrosynthetic analysis are passed to
this function.

\begin{algorithm}[H]
\caption{\textsc{ScoreVariant}$(s, \sigma, q, \textit{espsim})$}
\label{alg:score-variant}
\begin{algorithmic}[1]
\Require SMILES string $s$; specification $\sigma$;
         baseline quality $q \in [0,1]$;
         boolean \textit{espsim}
\Ensure  Feature dictionary, or $\bot$ for invalid inputs
\If{$s$ contains \texttt{`.'}}\Comment{multi-fragment: salt or mixture}
  \Return $\bot$
\EndIf
\State $m \leftarrow \textsc{ParseSMILES}(s)$
\If{$m = \bot$} \Return $\bot$ \EndIf
\State $\tau \leftarrow T(m, \sigma)$
       \Comment{Algorithm~\ref{alg:tanimoto}}
\State $\psi,\, (s_d)_{d\in\Desc} \leftarrow \textsc{Desirability}(m, \sigma)$
       \Comment{Algorithm~\ref{alg:desirability}}
\State $e \leftarrow \textsc{ESPsim}(m, \sigma)$ \textbf{if} \textit{espsim}
       \textbf{else} $\bot$
       \Comment{Algorithm~\ref{alg:espsim}}
\State $b \leftarrow \textsc{AssignBin}(q)$
       \Comment{Definition~\ref{def:bin}}
\Return $\bigl\{$
  \texttt{tanimoto}: $\tau$,\;
  \texttt{desirability}: $\psi$,\;
  $(s_d)_{d\in\Desc}$,\;
  \texttt{espsim}: $e$,\;
  \texttt{quality\_bin}: $b$
$\bigr\}$
\end{algorithmic}
\end{algorithm}


% ════════════════════════════════════════════════════════════════════════════
\section{Hit Classification}
\label{sec:hit}
% ════════════════════════════════════════════════════════════════════════════

\begin{definition}[Hit indicator]
\label{def:hit}
Given thresholds $\tau_t \in (0,1]$ and $\tau_d \in (0,1]$, variant $m$
is a \emph{hit} with respect to specification $\sigma$ if and only if
\[
  h(m, \sigma;\, \tau_t, \tau_d)
  \;=\; \1{T(m, \sigma) \geq \tau_t}
        \cdot
        \1{\Psi(m, \sigma) \geq \tau_d}
  \;=\; 1.
\]
The hit requires simultaneous satisfaction of the substructural and
physicochemical conservation gates; neither criterion can substitute for
the other.
\end{definition}

\begin{definition}[Hit set]
\label{def:hitset}
Let $\Var_M(\sigma)$ denote the set of synthesizable variants produced by
method $M$ for specification $\sigma$ (after Gate~1 filtering).
The \emph{hit set} of method $M$ at thresholds $(\tau_t, \tau_d)$ is
\[
  \Hit_M(\tau_t, \tau_d)
  \;=\; \bigl\{\, m \in \Var_M(\sigma)
               : h(m, \sigma;\, \tau_t, \tau_d) = 1 \,\bigr\}.
\]
\end{definition}

\begin{remark}
Results are reported at all combinations in
$\{0.6, 0.7, 0.85\} \times \{0.8\}$ for $(\tau_t, \tau_d)$.
No conclusion is drawn from a single threshold value.
\end{remark}


% ════════════════════════════════════════════════════════════════════════════
\section{Expansion Factor}
\label{sec:expansion}
% ════════════════════════════════════════════════════════════════════════════

\begin{definition}[Expansion factor]
\label{def:expansion}
Let $B$ denote the repeated-PrexSyn-sampling baseline.
The \emph{expansion factor} of method $M$ at thresholds $(\tau_t, \tau_d)$ is
\[
  \text{EF}(M;\, \tau_t, \tau_d)
  \;=\; \frac{|\Hit_M(\tau_t, \tau_d)|}{|\Hit_B(\tau_t, \tau_d)|}.
\]
If $|\Hit_B| = 0$, we define $\text{EF}(M) = +\infty$.
\end{definition}

The expansion factor is the \emph{primary ranking metric}: a value greater
than~1 indicates that method $M$ produces synthesizable, property-conserving
candidates that repeated PrexSyn sampling cannot reach.

\begin{algorithm}[H]
\caption{\textsc{ExpansionFactor}$(\Hit_M, \Hit_B)$}
\label{alg:expansion}
\begin{algorithmic}[1]
\Require Hit sets $\Hit_M$, $\Hit_B$ (sets of canonical SMILES strings)
\Ensure  $\text{EF} \in [0, +\infty)$
\If{$|\Hit_B| = 0$} \Return $+\infty$ \EndIf
\Return $|\Hit_M| \;/\; |\Hit_B|$
\end{algorithmic}
\end{algorithm}


% ════════════════════════════════════════════════════════════════════════════
\section{Chemical Diversity}
\label{sec:diversity}
% ════════════════════════════════════════════════════════════════════════════

\begin{definition}[Chemical diversity]
\label{def:diversity}
The \emph{chemical diversity} of a set of molecules $\mathcal{H}$ is
\[
  \text{Div}(\mathcal{H})
  \;=\; 1 - \frac{2}{|\mathcal{H}|\,(|\mathcal{H}|-1)}
            \sum_{\substack{m,\, m' \in \mathcal{H} \\ m \neq m'}}
            T\!\bigl(\fp{m},\, \fp{m'}\bigr)
  \;\in\; [0, 1].
\]
A value of~1 indicates that no two hits share any circular substructure
features at ECFP4 radius~2; a value of~0 indicates that all hits are
identical.
\end{definition}

\begin{remark}
$\text{Div}(\mathcal{H})$ is undefined for $|\mathcal{H}| < 2$;
the implementation returns 0 in this case.
\end{remark}


% ════════════════════════════════════════════════════════════════════════════
\section{Complementarity Analysis}
\label{sec:complementarity}
% ════════════════════════════════════════════════════════════════════════════

\subsection{Pairwise Jaccard overlap}

\begin{definition}[Pairwise Jaccard overlap]
\label{def:jaccard}
For methods $M_i, M_j \in \mathcal{M}$, the \emph{Jaccard overlap} of
their hit sets at fixed $(\tau_t, \tau_d)$ is
\[
  J(M_i, M_j)
  \;=\; \frac{|\Hit_{M_i} \cap \Hit_{M_j}|}{|\Hit_{M_i} \cup \Hit_{M_j}|}.
\]
$J = 0$ indicates fully complementary methods (disjoint hit sets);
$J = 1$ indicates fully redundant methods (identical hit sets).
\end{definition}

\begin{algorithm}[H]
\caption{\textsc{Complementarity}$\!\bigl((\Hit_M)_{M \in \mathcal{M}}\bigr)$
         --- pairwise Jaccard matrix}
\label{alg:complementarity}
\begin{algorithmic}[1]
\Require Hit sets $\Hit_M$ for each $M \in \Methods$
\Ensure  Matrix $J \in [0,1]^{|\Methods| \times |\Methods|}$
\For{$M_i \in \Methods$}
  \For{$M_j \in \Methods$}
    \State $u \leftarrow |\Hit_{M_i} \cup \Hit_{M_j}|$
    \If{$u = 0$}
      \State $J_{ij} \leftarrow 0$
    \Else
      \State $J_{ij} \leftarrow |\Hit_{M_i} \cap \Hit_{M_j}| \;/\; u$
    \EndIf
  \EndFor
\EndFor
\Return $J$
\end{algorithmic}
\end{algorithm}

\subsection{Optimal combination search}

\begin{definition}[Optimal $k$-method combination]
\label{def:best-combo}
For $k \in \{2, 3\}$, the \emph{optimal $k$-method combination} is
\[
  C^*(k)
  \;=\; \argmax_{\substack{C \subseteq \Methods \setminus \{B\} \\ |C| = k}}
        \left|\bigcup_{M \in C} \Hit_M(\tau_t, \tau_d)\right|.
\]
The baseline $B$ is excluded from the search; results are reported
for both values of $k$.
\end{definition}

\begin{algorithm}[H]
\caption{\textsc{BestCombinations}$\!\bigl((\Hit_M)_{M \in \mathcal{M}},\, k_\text{top}\bigr)$}
\label{alg:best-combos}
\begin{algorithmic}[1]
\Require Hit sets; number of results to return $k_\text{top}$
\Ensure  Ranked list of (combination, coverage) pairs
\State $\text{results} \leftarrow [\;]$
\For{$k \in \{2, 3\}$}
  \For{each $C \subseteq \Methods \setminus \{B\}$ with $|C| = k$}
    \State $\text{coverage} \leftarrow \left|\bigcup_{M \in C} \Hit_M\right|$
    \State $\text{results.append}(C, \text{coverage})$
  \EndFor
\EndFor
\Return top $k_\text{top}$ entries of results sorted by coverage descending
\end{algorithmic}
\end{algorithm}


% ════════════════════════════════════════════════════════════════════════════
\section{Stratification and Multi-Threshold Summary}
\label{sec:summary}
% ════════════════════════════════════════════════════════════════════════════

\begin{definition}[Baseline quality and stratification bins]
\label{def:bin}
The \emph{baseline quality} of a specification $\sigma$ is the Tanimoto
similarity of the PrexSyn seed candidate to the spec fingerprint:
\[
  q(\sigma) \;=\; T(m_\text{seed},\, \sigma),
  \qquad m_\text{seed} = \argmax_{m \in \Var_\text{PrexSyn}} T(m, \sigma).
\]
Specifications are assigned to one of four \emph{difficulty bins}:
\[
  b(\sigma) \;=\;
  \begin{cases}
    B_1 & q(\sigma) < 0.5 \\
    B_2 & 0.5 \leq q(\sigma) < 0.7 \\
    B_3 & 0.7 \leq q(\sigma) < 0.85 \\
    B_4 & q(\sigma) \geq 0.85.
  \end{cases}
\]
\end{definition}

\begin{definition}[Per-method summary statistics]
\label{def:summary}
For each method $M$, threshold $\tau_t$, fixed $\tau_d = 0.8$, and
bin $B_i$, let $\mathcal{V}_{M,i}$ and $\Hit_{M,i}$ denote the variant
and hit sets restricted to specifications in bin $B_i$.
The summary statistics are:

\medskip
\begin{tabular}{ll}
  Hit rate & $\displaystyle r_{M,i}(\tau_t) = \frac{|\Hit_{M,i}(\tau_t)|}{|\mathcal{V}_{M,i}|}$ \\[1.5ex]
  Unique hits & $\displaystyle n_{M,i}(\tau_t) = |\Hit_{M,i}(\tau_t)|$ \\[1.5ex]
  Expansion factor & $\displaystyle \text{EF}_{M,i}(\tau_t) = \dfrac{n_{M,i}(\tau_t)}{n_{B,i}(\tau_t)}$ \\[1.5ex]
  Mean Tanimoto & $\displaystyle \bar{T}_{M,i} = \dfrac{1}{|\mathcal{V}_{M,i}|} \sum_{m \in \mathcal{V}_{M,i}} T(m,\sigma)$ \\[1.5ex]
  Mean desirability & $\displaystyle \bar{\Psi}_{M,i} = \dfrac{1}{|\mathcal{V}_{M,i}|} \sum_{m \in \mathcal{V}_{M,i}} \Psi(m,\sigma)$ \\[1.5ex]
  Diversity & $\displaystyle \text{Div}_{M,i}(\tau_t) = \text{Div}\!\bigl(\Hit_{M,i}(\tau_t)\bigr)$
\end{tabular}
\end{definition}

\begin{algorithm}[H]
\caption{\textsc{Summarize}$(D,\, \boldsymbol{\tau_t},\, \tau_d)$ --- multi-threshold stratified summary}
\label{alg:summarize}
\begin{algorithmic}[1]
\Require Score DataFrame $D$ with columns
  $\{\texttt{method}, \texttt{quality\_bin}, \texttt{tanimoto},
     \texttt{desirability}, \texttt{variant\_smiles}\}$;
  threshold list $\boldsymbol{\tau_t} = \{0.6, 0.7, 0.85\}$;
  fixed $\tau_d = 0.8$
\Ensure  Summary DataFrame indexed by $(\text{method}, b, \tau_t)$
\State $\text{rows} \leftarrow [\;]$
\For{$\tau_t \in \boldsymbol{\tau_t}$}
  \State Compute $\Hit_B^{(b)}(\tau_t,\tau_d)$ for each bin $b$
         \Comment{baseline hit sets for EF}
  \For{each group $(M, b)$ in $D$}
    \State $\mathcal{V} \leftarrow $ rows in group; $\Hit \leftarrow $ rows where $h = 1$
    \State Compute $r, n, \text{EF}, \bar{T}, \bar{\Psi}, \text{Div}$
           per Definition~\ref{def:summary}
    \State $\text{rows.append}(M, b, \tau_t, r, n, \text{EF}, \bar{T}, \bar{\Psi}, \text{Div})$
  \EndFor
\EndFor
\Return $\text{DataFrame}(\text{rows})$ sorted by $(\tau_t, b, M)$
\end{algorithmic}
\end{algorithm}


% ════════════════════════════════════════════════════════════════════════════
\section*{References}
% ════════════════════════════════════════════════════════════════════════════

\begin{thebibliography}{9}

\bibitem{rdkit}
G.~Landrum et al.
\textit{RDKit: Open-source cheminformatics}, 2024.
\url{http://www.rdkit.org}

\bibitem{espsim}
G.~Bolcato, E.~Heid, S.~Riniker.
On the value of using 3D shape and electrostatic similarities in deep
generative methods.
\textit{J.\ Chem.\ Inf.\ Model.}\ \textbf{62}(6), 1388--1398 (2022).

\bibitem{morgan}
H.~L.~Morgan.
The generation of a unique machine description for chemical structures ---
a technique developed at Chemical Abstracts Service.
\textit{J.\ Chem.\ Doc.}\ \textbf{5}(2), 107--113 (1965).

\end{thebibliography}

\end{document}
